%% kapa: exam=1 round=33 year=2022 period=1 form=A subject=civil_law subject_ko=민법 seq=6 number=6
%% source: Q-Net 감정평가사 기출문제 (cst00309) artlSeq=5209997 file=raw/1008/5209997/2022년도 제33회 감정평가사 1차시험 1교시(민법,경제학,부동산학).pdf
%% conversion: kordoc 4.24.51 -> tools/gen1.py (auto); review=proofread-sonnet; flags=none
\begin{question}{6}
甲사단법인의 대표이사 乙이 외관상 그 직무에 관한 행위로 丙에게 불법행위를 한 경우에 관한 설명으로 옳지 않은 것은? (다툼이 있으면 판례에 따름)
\begin{choices}
  \item 乙의 불법행위로 인해 甲이 丙에 대해 손해배상책임을 지는 경우에도 乙은 丙에 대한 자기의 손해배상책임을 면하지 못한다.
  \item 甲의 손해배상책임 원인이 乙의 고의적인 불법행위인 경우에는 丙에게 과실이 있더라도 과실상계의 법리가 적용될 수 없다.
  \item 丙이 乙의 행위가 실제로는 직무에 관한 행위에 해당하지 않는다는 사실을 알았거나 중대한 과실로 알지 못한 경우에는 甲에게 손해배상책임을 물을 수 없다.
  \item 甲의 사원 丁이 乙의 불법행위에 가담한 경우, 丁도 乙과 연대하여 丙에 대하여 손해배상책임을 진다.
  \item 甲이 비법인사단인 경우라 하더라도 甲은 乙의 불법행위로 인한 丙의 손해를 배상할 책임이 있다.
\end{choices}
\end{question}
\answer{2}
